\documentclass[11pt,a4paper]{article}

\usepackage{kotex}
\usepackage{fontspec}
\setmainfont{Times New Roman}
\setmainhangulfont{AppleMyungjo}
\setsanshangulfont{Apple SD Gothic Neo}

\usepackage[a4paper,top=18mm,bottom=18mm,left=20mm,right=20mm]{geometry}
\usepackage{fancyhdr}
\setlength{\headheight}{24pt}
\usepackage{enumitem}
\usepackage{mdframed}
\usepackage{array}
\usepackage{booktabs}

\setlength{\parindent}{1em}
\linespread{1.18}

\pagestyle{fancy}
\fancyhf{}
\renewcommand{\headrulewidth}{0.6pt}
\fancyhead[L]{\sffamily\bfseries 2026 고1 영어 변형 — 지문 20}
\fancyhead[R]{\sffamily 자녀 가치 교육}
\renewcommand{\footrulewidth}{0pt}
\fancyfoot[C]{\framebox[16mm][c]{\thepage}}

\newcommand{\qno}[1]{\par\vspace{8pt}\noindent{\sffamily\bfseries #1.}\ }
\newcommand{\instr}[1]{{\sffamily #1}\par\vspace{3pt}}
\newlist{choices}{enumerate}{1}
\setlist[choices]{label=,leftmargin=1.5em,itemsep=2pt,topsep=3pt,parsep=0pt}
\newmdenv[linewidth=0.6pt,roundcorner=3pt,innertopmargin=7pt,innerbottommargin=7pt,
          innerleftmargin=9pt,innerrightmargin=9pt,skipabove=5pt,skipbelow=5pt]{pbox}
\newcommand{\nemo}[1]{\fbox{\,#1\,}}

\begin{document}

\begin{center}
{\fontsize{15}{17}\selectfont\sffamily\bfseries [지문 20] 다음 글을 읽고 물음에 답하시오.}
\end{center}
\vspace{4pt}

\begin{pbox}
Remember how you learned to drive or cook? You practiced while someone coached you,
reminding you what to do until you were able to coach yourself and then, eventually, do it
automatically. Children learn values much the same way. They practice different kinds of
behavior, while, you, as coach, help focus their attention on what is important and on
\rule{3cm}{0.4pt} important skills. You support them with your praise, encouragement and
gentle reminders. If you don't coach your child, she will find her coaches elsewhere and be
guided by the values of the media, her peers and anyone else who captures her interest.
So, step up to the plate, don't be afraid and help your child learn how to be a good person,
step by step.
\end{pbox}

\qno{1}\instr{다음 빈칸에 들어갈 말로 가장 적절한 것은?}
\begin{choices}
\item ① memorizing
\item ② refining
\item ③ ignoring
\item ④ replacing
\item ⑤ evaluating
\end{choices}

\qno{2}\instr{윗글의 제목으로 가장 적절한 것은?}
\begin{choices}
\item ① Why Children Learn Better from Peers Than from Parents
\item ② The Role of Media in Shaping Children's Behavior
\item ③ How Automatic Skills Develop Through Repetition
\item ④ Coach Your Child: Be the Guide They Need to Grow Well
\item ⑤ Encouraging Independence: When to Stop Coaching Your Child
\end{choices}

\qno{3}\instr{[서술형] 다음 문장을 우리말로 자연스럽게 해석하시오.}

\begin{quote}
\itshape
If you don't coach your child, she will find her coaches elsewhere and be guided by the
values of the media, her peers and anyone else who captures her interest.
\end{quote}

\noindent $\Rightarrow$ \rule{\linewidth}{0.4pt}

\vspace{4pt}
\noindent \rule{\linewidth}{0.4pt}

% ===================== Q4 =====================
\qno{4}\instr{윗글의 내용을 가장 잘 요약한 것은?}
\begin{choices}
\item ① Children should be left alone so that they can discover their own values naturally.
\item ② Parents should coach children in values before outside influences become their guides.
\item ③ Media and peers are more effective than parents at teaching social behavior.
\item ④ Children learn values only after they have mastered practical skills like driving or cooking.
\item ⑤ Praise should be avoided because it prevents children from becoming independent.
\end{choices}

\end{document}
